%%
%% This is a placeholder manuscript shell for IMC'27 (ACM Internet Measurement Conference).
%% Status: PENDING (Sequential Plan B fallback; activate only after explicit WWW'27 rejection).
%%
\documentclass[sigconf,anonymous,review]{acmart}

\acmConference[IMC '27]{ACM Internet Measurement Conference 2027}{October 2027}{TBD}
\acmYear{2027}
\acmDOI{}
\acmISBN{}

\title{Measuring the Model Context Protocol: Remote Endpoints, Secret Headers, and Domain Migrations in the Agentic Web}

\author{Anonymous Author(s)}
\affiliation{%
  \institution{Anonymous Institution}
  \country{Anonymous Country}
}
\email{anonymous@example.org}

\begin{abstract}
The Model Context Protocol (MCP) establishes a standardised web application interface connecting autonomous language model agents to remote microservices. We perform a comprehensive, longitudinal measurement of the official MCP registry across 126,847 version records and 37,815 distinct servers. Our network measurement reveals that 61.89\% of registered servers declare remote HTTP endpoints and 56.96\% operate purely as remote web services without installable packages. These endpoints resolve to 15,738 distinct hosts across 6,432 registrable domains, with 79.5\% unranked in the Tranco Top-1M. Furthermore, 3,679 endpoints demand confidential headers marked \texttt{isSecret}, binding ambient credentials to outbound web requests. Across 89,032 publication-ordered version transitions, releases arrive rapidly (median interval 0.8 days, 52.4\% within 24 hours). Longitudinal tracking reveals 1,353 transitions modifying remote URLs, with 52.4\% changing destination hostnames and 36.1\% crossing administrative eTLD+1 boundaries, predominantly migrating to multi-tenant serverless or PaaS platforms. Syntactic pattern screening across 40,672 tool descriptions yields a 1.25\% flag rate, but recovers only 6.2\% of benchmark prompt injection attacks while flagging benign documentation in 89.2\% of real-world hits. We present architectural recommendations for origin-bound credential scoping, out-of-band manifest attestations, and registry domain lifecycle verification to secure agentic web infrastructure.
\end{abstract}

\begin{CCSXML}
<ccs2012>
   <concept>
       <concept_id>10003033.10003079</concept_id>
       <concept_desc>Networks~Network performance evaluation</concept_desc>
       <concept_significance>500</concept_significance>
   </concept>
   <concept>
       <concept_id>10002951.10003260.10003282</concept_id>
       <concept_desc>Information systems~Web applications</concept_desc>
       <concept_significance>500</concept_significance>
   </concept>
</ccs2012>
\end{CCSXML}

\ccsdesc[500]{Networks~Network performance evaluation}
\ccsdesc[500]{Information systems~Web applications}

\keywords{Internet measurement, Model Context Protocol, agentic systems, HTTP endpoints, domain migrations, Web infrastructure}

\begin{document}

\maketitle

\section{Introduction}
The rapid proliferation of agentic AI systems has catalysed standardised protocols for connecting foundation models to external web tools and APIs. The Model Context Protocol (MCP) has emerged as an open standard for this interface. While often characterised as a local client extension mechanism, our measurement reveals that MCP functions primarily as a distributed Web protocol, with the majority of servers exposing remote HTTP streaming endpoints.

\section{Measurement Methodology and Dataset}
We analyze a longitudinal census of the official registry walk covering 126,847 records and 37,815 servers, evaluating remote endpoints, Tranco domain authority rankings, and secret-header configurations.

\section{Longitudinal URL and Domain Migrations}
We identify 1,353 URL change events across consecutive version releases, categorising hostname migrations, cross-eTLD+1 shifts, and transitions into serverless multi-tenant hosting environments.

\section{Discussion and Protocol Recommendations}
We discuss operational risks associated with unverified remote HTTP tool endpoints and recommend origin-bound credential scoping and manifest attestation architectures.

\bibliographystyle{ACM-Reference-Format}
\bibliography{references}

\end{document}
